\documentclass[10pt,a4paper]{amsart}
\usepackage[margin=19mm]{geometry}
\usepackage{amsmath,amssymb,mathtools}
\usepackage[scr=rsfs,cal=euler]{mathalfa}
\usepackage{xfrac}
\usepackage[unicode,hidelinks,bookmarks=false]{hyperref}
\newcommand{\ie}{i.e.\ }
\newcommand{\cf}{cf.\ }
\newcommand{\resp}{respectively\ }
\newcommand{\Subsection}{\subsection}
\providecommand{\nobreakdash}{\nobreak\hbox{-}}
\setlength{\emergencystretch}{2em}
\multlinegap0pt
\usepackage{bm}
\usepackage{bbm}
\usepackage{dsfont}
\usepackage{xspace}
\usepackage{cancel}
\usepackage{mathtools}
\usepackage{amsthm}
\makeatletter
\@ifundefined{thm}{\newtheorem{thm}{Thm}}{}
\@ifundefined{prop}{\newtheorem{prop}[thm]{Proposition}}{}
\makeatother
\makeatletter
\expandafter\ifx\csname proof\endcsname\relax
\newenvironment{proof}{\noindent\emph{Proof.}}{\hfill$\square$\medskip}
\fi
\expandafter\ifx\csname rmk\endcsname\relax
\newenvironment{rmk}{\medskip\noindent\emph{Remark.}}{\medskip}
\fi
\newcommand{\sph}{\mathbb{S}}
\makeatother
\title[One-dimensional half-harmonic maps into the circle and their]{One-dimensional half-harmonic maps into the circle and their degree}
\author{Luca Martinazzi, Ali Hyder}
\date{Source: arXiv:2405.18037v2 (CC BY 4.0)}
\begin{document}
\maketitle

\noindent\emph{Source and licence.} One-dimensional half-harmonic maps into the circle and their degree. Version arXiv:2405.18037v2 (CC BY 4.0), distributed under the Creative Commons Attribution 4.0 International licence, as recorded in the arXiv licence metadata for this exact version and shown on the abstract page https://arxiv.org/abs/2405.18037v2. Full rights notice: licenses/arxiv-2405.18037-CC-BY-4.0.txt.\par\smallskip

\noindent\textit{Editorial scope.} Counted unit: the Prop whose source label is \texttt{p:degjump} in the article (source0.tex lines 443--447, source label \texttt{p:degjump}), delivered together with the article's own complete original proof of it (source0.tex lines 449--566, one \texttt{proof} environment of 118 lines). No mathematics is written, repaired or re-derived: statement and proof are the article's own text line for line. Required context retained uncounted: degformula2 (lines 309--311); defw (lines 420--426); defuv (lines 428--434); trace (lines 1221--1223). Every definition, display and label of those lines is retained unchanged. Omitted from this package and not redistributed: 1--308, 312--419, 427--427, 435--442, 448--448, 567--1220, 1224--1496. Nothing inside a retained excerpt is omitted or altered: every retained body is delivered exactly as the article's lines are written, so no omission receipt of any kind accompanies this package. Citations: the retained excerpts carry no citation command at all, so no bibliography entry of the article is transcribed and no reference is redistributed. Reference adaptation: none was needed and none was made; every reference inside the delivered unit resolves inside it, so no unresolved reference or citation is produced. Printed slips kept literal: no printed word, symbol, hypothesis or number of the counted statement or of its proof was altered. Layout and class: the article's own class is \texttt{article} with packages including amsfonts, amssymb, amsmath, amscd, latexsym, makeidx, theorem, hyperref, color, txfonts; the wrapper is the lane's pinned \texttt{amsart} editorial wrapper at 10pt on A4, which restates the theorem-like environments the retained text needs and adds the article's own macro definitions that the retained text needs (\texttt{\textbackslash sph}), with extra wrapper packages (bm, bbm, dsfont, xspace, cancel, mathtools). The delivered numbering is the unit's own. Assets: no retained span contains a figure environment, a graphics-inclusion command, a PDF-page inclusion, a table or a TikZ picture; no image file is copied into this package, the package declares no asset dependency and no image file is redistributed with it. Licence: version arXiv:2405.18037v2 of this article is distributed under the Creative Commons Attribution 4.0 International licence, as recorded in the arXiv licence metadata for this exact version and shown on the abstract page https://arxiv.org/abs/2405.18037v2. A full rights notice is delivered at licenses/arxiv-2405.18037-CC-BY-4.0.txt. Characters: every retained excerpt is pure ASCII, so the delivered engine, XeLaTeX with its default Latin Modern text font, reports no missing character.
\begin{equation}\label{degformula2}
\deg(u):=\frac{1}{\pi}\int_D JU dxdy,\quad u\in H^{\frac{1}{2},2}(\sph^1,\sph^1), \, U\in H^{1,2}(D^2,\mathbb{C}), \, U|_{\partial D^2}=u.
\end{equation}

\begin{equation}\label{defw}
v\& u(e^{i\theta}):=
\begin{cases}
v(e^{i\theta}) &\text{for }0\le \theta\le \pi\\
u(e^{i\theta}) &\text{for } \pi <\theta< 2\pi,
\end{cases}
\end{equation}

\begin{equation}\label{defuv}
v\# u(e^{i\theta}):=
\begin{cases}
v(e^{i\theta}) &\text{for }0\le \theta\le \pi\\
u(e^{-i\theta}) &\text{for } \pi <\theta< 2\pi.
\end{cases}
\end{equation}

\begin{prop}\label{trace}
For any bounded Lipschitz domain $\Omega$ with $\Gamma:=\partial\Omega$, the trace maps $u\mapsto u|_\Gamma$ is surjective from $H^{1,2}(\Omega)$ to $H^{\frac12,2}(\Gamma)$.
\end{prop}

\begin{prop}\label{p:degjump} Consider $u \in H^{\frac{1}{2},2}(\sph^1,\sph^1)$ and $v:\sph^1_+\to \sph^1$. Set $v\&u$ and $v\#u$ as in \eqref{defw} and \eqref{defuv} and assume that $v\&u, v\#u\in H^{\frac{1}{2},2}(\sph^1,\sph^1)$. Then
\begin{equation}\label{eq:jump}
\deg(v\&u)= \deg(u)+\deg (v\#u).
\end{equation}
\end{prop}

\begin{proof}
Consider $\gamma:[-1,1]\times\{0\}\to\sph^1_-$ given by $\gamma(x,0)=e^{i\frac{\pi}{2}(x-1)}$. Clearly $\gamma$ is a diffeomorphism and
\begin{equation}\label{diff}
|\gamma'|\equiv\frac{\pi}{2},\quad |(\gamma^{-1})'|\equiv\frac{2}{\pi}.
\end{equation}
On $[-1,1]\times \{0\}$ define the function
$$u^*(x,0)=u(\gamma(x,0))=u(e^{i\frac{\pi}{2}(x-1)}).$$
Now define the maps $u^*_\pm :\partial D^2_\pm \to \sph^1$
\[
\begin{cases}
u^*_\pm :=u&\text{on }\sph^1_\pm\\
u^*_\pm =u^*&\text{on } (-1,1)\times \{0\},
\end{cases}
\]
and $v^*:\partial D^2_+\to \sph^1$
\[
\begin{cases}
v^*=v&\text{on }\sph^1_+\\
v^*=u^*&\text{on } (-1,1)\times \{0\}.
\end{cases}
\]
Since  $|(x,0)-(\xi,0)| \ge \frac{2}{\pi}|\gamma(x,0)-\gamma(\xi,0)|$   for $x, \xi\in [-1,1]$, also using \eqref{diff}, we get 
\begin{equation}\label{stimau*}
\begin{split}
\int_{-1}^{1}\int_{-1}^{1}\frac{|u^*(x,0)-u^*(\xi,0)|^2}{|(x,0)-(\xi,0)|^2}dxd\xi
&\le \frac{\pi^2}{4} \int_{-1}^{1}\int_{-1}^{1}\frac{|u(\gamma(x,0))-u(\gamma(\xi,0))|^2}{|\gamma(x,0)-\gamma(\xi,0)|^2}dxd\xi\\
&=  \int_{\sph^1_-}\int_{\sph^1_-}\frac{|u(e^{is})-u(e^{it})|^2}{|e^{is}-e^{it}|^2}dsdt.
\end{split}
\end{equation}
We now claim
\begin{align}\label{claimgamma0}
|e^{is}-(x,0)|&\ge \frac{|e^{is}-\gamma(x,0)|}{2},\quad \text{for }e^{is}\in \sph^1_-, \, x\in [-1,1],\\
\label{claimgamma}
|e^{is}-(x,0)|&\ge \frac{|e^{is}-\gamma(x,0)|}{4},\quad \text{for }e^{is}\in \sph^1_+,\ x\in [-1,1].
\end{align}
%Since
%$$|e^{is}-(x,0)|= |e^{-is}-(x,0)|, \quad |e^{is}-\gamma(x,0)|\ ge |e^{-is}-\gamma(x,0)|,\quad \text{for }e^{is}\in \sph^1_+,$$
To prove \eqref{claimgamma0}, by symmetry it suffices to consider the case $-\frac{\pi}{2}\le s\le 0$. We consider $2$ cases.

\noindent\textbf{Case 1:} $0\le \frac{\pi}{4}(1-x)\le -s\le\frac{\pi}{2}$. 
We have
\begin{equation}\label{eqgamma}
|e^{is}-(x,0)|\ge |\sin s| \ge \frac{2}{\pi}|s|\ge\frac{|s|}{2}
\end{equation}
(the second inequality follows from showing that $f(t)=\frac{\sin t}{t}$ is decreasing for $0<t\le \frac{\pi}{2}$).
Moreover, estimating $|e^{it_1}-e^{it_2}|\le |t_1-t_2|$, we get
$$|e^{is}-\gamma(x,0)|\le \left|s+\frac{\pi}{2}(1-x)\right|\le |s|\le 2  |e^{is}-(x,0)|.$$
\noindent\textbf{Case 2:} $-s\le \frac{\pi}{4}(1-x)\le \frac{\pi}{2}$. We have
\begin{equation}\label{eqgamma2}
|e^{is}-(x,0)|\ge 1-x,
\end{equation}
since the circle centred at $(x,0)$ of radius $1-x$ is contained $D^2$ and $e^{is}\in \partial D^2$. Then
$$|e^{is}-\gamma(x,0)|\le \left|s+\frac{\pi}{2}(1-x)\right|\le \frac{\pi}{2}(1-x) \le  2|e^{is}-(x,0)|.$$
Then \eqref{claimgamma0} is proven. In order to prove \eqref{claimgamma}, again by symmetry, it suffices to prove it for $0\le s\le \frac{\pi}{2}$ and we consider three cases.

%\noindent \textbf{Case 1}: $\frac{\pi}{6}\le s\le \frac{\pi}{2}$. In this case
%$$|e^{is}-(x,0)|\ge |\sin (s)|\ge \frac{1}{2} \ge \frac{|e^{is}-\gamma(x,0)|}{4}.$$

\noindent\textbf{Case 1:} $0\le x\le 1$, $ \frac{\pi}{2}(1-x)\le s\le \frac{\pi}{2}$.
As before we have \eqref{eqgamma}.
Moreover $|e^{is}-\gamma(x,0)|\le s+\frac{\pi}{2}(1-x)$.  Hence, by our assumption,
$$\frac{|e^{is}-\gamma(x,0)|}{4}\le \frac{s}{4}+ \frac{\pi}{8}(1-x)\le \frac{s}{4}+\frac{s}{4}\le |e^{is}-(x,0)|.$$

\noindent\textbf{Case 2:} $0\le x\le 1$, $0\le s\le \frac{\pi}{2}(1-x)\le \frac{\pi}{2}$. We have \eqref{eqgamma2} as before. Then
$$|e^{is}-\gamma(x,0)|\le s+\frac{\pi}{2}(1-x)\le \pi (1-x)\le 4(1-x)\le 4 |e^{is}-(x,0)|.$$

\noindent\textbf{Case 3:} $-1\le  x< 0$, $0\le s\le \frac{\pi}{2}$. Since
$|e^{is}-(x,0)|\ge 1$, we have
$$|e^{is}-\gamma(x,0)|\le 2 \le  2 |e^{is}-(x,0)|. $$
Then also \eqref{claimgamma} is proven.

\medskip


Now \eqref{claimgamma0} together with \eqref{diff} gives
\[\begin{split}
\int_{[-1,1]\times \{0\}}\int_{\sph^1_-} \frac{|u^*_-(e^{is})-u^*_-(x,0)|^2}{|e^{is}-(x,0)|^2}ds dx&\le {4} \int_{[-1,1]\times \{0\}}\int_{\sph^1_-} \frac{|u(e^{is})-u(\gamma(x,0))|^2}{|e^{is}-\gamma(x,0)|^2}dsdx\\
&= \frac{8}{\pi} \int_{\sph^1_-}\int_{\sph^1_-}\frac{|u(e^{is})-u(e^{it})|^2}{|e^{is}-e^{it}|^2}dsdt.
\end{split}\]
With \eqref{stimau*}, it follows that
$$[u^*_-]_{H^{\frac12,2}(\partial D^2_-)}\le C[u]_{H^{\frac12,2}(\sph^1_-)}.$$
For \eqref{claimgamma} we infer
\[\begin{split}
\int_{[-1,1]\times \{0\}}\int_{\sph^1_+} \frac{|u^*_+(e^{is})-u^*_+(x,0)|^2}{|e^{is}-(x,0)|^2}dsdx& \le 16\int_{[-1,1]\times \{0\}}\int_{\sph^1_+} \frac{|u(e^{is})-u(\gamma(x,0))|^2}{|e^{is}-\gamma(x,0)|^2}dsdx\\
&=\frac{32}{\pi}  \int_{\sph^1_-}\int_{\sph^1_+}\frac{|u(e^{is})-u(e^{it})|^2}{|e^{is}-e^{it}|^2}dsdt.
\end{split}\]
Similarly
$$\int_{[-1,1]\times \{0\}}\int_{\sph^1_+} \frac{|v^*(e^{is})-v^*(x,0)|^2}{|e^{is}-(x,0)|^2}dsdx\le \frac{32}{\pi} \int_{\sph^1_-}\int_{\sph^1_+}\frac{|v(e^{is})-u(e^{it})|^2}{|e^{is}-e^{it}|^2}dsdt.$$
Together with \eqref{stimau*}, it follows that 
$$[u^*_+]_{H^{\frac12,2}(\partial D^2_+)}\le C[u]_{H^{\frac12,2}(\sph^1)},\qquad [v^*]_{H^{\frac12,2}(\partial D^2_+)}\le C[v\&u]_{H^{\frac12,2}(\sph^1)}.$$
Let us now call $U_+\in H^{1,2}(D^2_+,\mathbb{C})$, $U_-\in H^{1,2}(D^2_-,\mathbb{C})$,  $V\in H^{1,2}(D^2_+,\mathbb{C})$ any extensions of $u^*_+$, $u^*_-$ and $v^*$, respectively, which exist thanks to Proposition \ref{trace}.
Since $U_+$, $U_-$ and $V$ have the same trace $u^*$ on $(-1,1)\times \{0\}$ they can be patched together to obtain extensions $U$, $V\&U$ and $V\#U$ of $u$, $v\&u$ and  $v\#u$, respectively, namely
\[
U=\begin{cases}
U_+&\text{in }D^2_+\\
U_-&\text{in }D^2_-,
\end{cases}
\qquad
V\&U=\begin{cases}
V&\text{in }D^2_+\\
U_-&\text{in }D^2_-,
\end{cases}
\]
and
\[
V\#U(z)=\begin{cases}
V(z)&\text{for }z\in D^2_+\\
U_+(\bar z)&\text{for }z\in D^2_-
\end{cases}
\]
Then, thanks to \eqref{degformula2}
\[\begin{split}
\deg(u\&v)&=\frac{1}{\pi}\int_{D^2} J(V\&U)\, dxdy\\
&=\frac{1}{\pi}\left(\int_{D^2_+} JV\, dxdy-\int_{D^2_+} JU_+\, dxdy\right)+ \frac{1}{\pi}\left(\int_{D^2_+} JU_+\, dxdy+\int_{D^2_-} JU_-\, dxdy\right)\\
&=\frac{1}{\pi}\int_{D^2} J(V\#U)\, dxdy+ \frac{1}{\pi}\int_{D^2} JU\, dxdy\\
&=\deg(v\#u)+\deg(u).
\end{split}\]
\end{proof}

\end{document}
